% Part: first-order-logic
% Chapter: axiomatic-deduction
% Section: axioms-rules-propositional

\documentclass[../../../include/open-logic-section]{subfiles}

\begin{document}

\iftag{FOL}
      {\olfileid[hi]{fol}{axd}{prp}}
      {\olfileid[hi]{pl}{axd}{prp}}

\olsection{प्रतिज्ञप्ति-संयोजकों के अभिगृहीत और नियम}

\begin{defn}[अभिगृहीत]
प्रतिज्ञप्ति-संयोजकों के \emph{अभिगृहीतों} का समुच्चय $\PAx$ निम्न रूपों के
सभी !!{formula}s से बना है:
\begin{align}
  & (!A \land !B) \lif !A \ollabel{ax:land1}\\
  & (!A \land !B) \lif !B \ollabel{ax:land2}\\
  & {}!A \lif (!B \lif (!A \land !B)) \ollabel{ax:land3}\\
  & {}!A \lif (!A \lor !B) \ollabel{ax:lor1}\\
  & {}!A \lif (!B \lor !A) \ollabel{ax:lor2}\\
  & (!A \lif !C) \lif ((!B \lif !C) \lif ((!A \lor !B) \lif !C)) \ollabel{ax:lor3}\\
  & {}!A \lif (!B \lif !A) \ollabel{ax:lif1}\\
  & (!A \lif (!B \lif !C)) \lif ((!A \lif !B) \lif (!A \lif !C)) \ollabel{ax:lif2}\\
  & (!A \lif !B) \lif ((!A \lif \lnot !B) \lif \lnot !A) \ollabel{ax:lnot1}\\
  & \lnot !A \lif (!A \lif !B) \ollabel{ax:lnot2}\\
  & \ltrue \ollabel{ax:ltrue}\\
  & \lfalse \lif !A \ollabel{ax:lfalse1}\\
  & (!A \lif \lfalse) \lif \lnot !A \ollabel{ax:lfalse2}\\
  & \lnot\lnot !A \lif !A \ollabel{ax:dne}
\end{align}
\end{defn}

\begin{defn}[पूर्वपक्ष-स्थापन]
  यदि $!B$ और $!B \lif !A$ पहले ही !!a{derivation} में आ चुके हों, तो
  $!A$ एक सही अनुमान-चरण है।
\end{defn}

पूर्वपक्ष-स्थापन नियम को संक्षेप में ``\MP'' लिखेंगे।

\end{document}
